% TOP_PROBLEM: {"id": 30003909, "title": "Axisymmetry of Spherical Near-Horizon Geometries", "category_id": 16, "category": {"id": 16, "name": "physics", "display_name": "Mathematical Physics", "description": "Problems at the intersection of mathematics and physics.", "slug": "physics", "order_index": 16, "created_at": "2026-07-31T15:26:25.671Z"}, "status": "partially_solved", "statusQualification": "The actual vacuum equation was recovered from PDF page 18. The report establishes local rigidity near Kerr and states the global removal-of-axisymmetry conjecture; modern status is not freshly established here. The problem is not enlarged to matter or nonzero cosmological constant."}
% FIRST_POSTED: 2026-10-11
% ENTRY_KIND: statement_only
% SUBMISSION: {"kind": "statement_only", "category": "top_problems", "problem_id": "30003909", "proposal_type": "established", "source": "Mathematical General Relativity (2018), Oberwolfach Reports, DOI 10.4171/owr/2018/36, pp. 2157--2251. Primary report accessed 2026-10-11; relevant PDF pages 18. https://doi.org/10.4171/owr/2018/36; Chrusciel--Szybka--Tod 2018 (catalogue-listed context; not independently consulted in this preparation). https://doi.org/10.1088/1361-6382/aa90e7", "status_evidence": "Canonical source status partially_solved retained. Supplied literature triage dated 2026-08-21: Kerr near-horizon uniqueness is proved in a neighborhood of the Kerr family, including the non-axisymmetric perturbative classification, but a global uniqueness theorem without the symmetry restriction is not established. Primary-report access checked 2026-10-11; this historical report and inherited triage do not certify current resolution of every component.", "status_checked_on": "2026-10-11", "duplicate_check": "Compared canonical ID, title and normalized target with the complete local UnsolvedMath registry; normalized exact-text matches: none; title matches: none. Related or overlapping canonical records: none identified in this batch. Primary source context and the existing OWR statement audit used for scope and duplicate review; an exhaustive mathematical-equivalence/literature search was not performed.", "attribution": "Original proposer not named in the supplied canonical record; attribution remains to the identified Oberwolfach contribution/report, not its editors as invented proposers. Explanatory statement preparation and source/semantic audit assisted by OpenAI GPT-6 Codex; no human mathematical authorship or independent proof verification is claimed."}
% SOURCE_STATEMENT: "Is every solution of the near-horizon-geometry equations on $S^2$ axisymmetric and therefore induced by a degenerate Kerr metric?"
% Preparation and final self-audit: CONTRIBUTING.md prompts C and D; source/audit artifacts retained in the local physics_statements_2026_10_11 batch.
% Imported from: physics_statements.tex; UnsolvedMath ID 30003909
% !TeX program = lualatex
\documentclass[11pt]{article}
\usepackage[margin=25mm]{geometry}
\usepackage{fontspec}
\setmainfont{Latin Modern Roman}
\usepackage{amsmath,amssymb,mathtools}
\usepackage{unicode-math}
\setmathfont{Latin Modern Math}
\usepackage{xurl,hyperref}
\hypersetup{colorlinks=true,urlcolor=blue}
\setcounter{secnumdepth}{0}
\setlength{\parindent}{0pt}
\setlength{\parskip}{5pt}
\setlength{\emergencystretch}{3em}
\newcommand{\sref}[2]{\hyperlink{source:#1:#2}{[S#2]}}
\newcommand{\cataloguescope}[1]{\paragraph{Recorded target.}#1}
\newcommand{\R}{\mathbb R}
\newcommand{\C}{\mathbb C}
\newcommand{\N}{\mathbb N}
\newcommand{\Z}{\mathbb Z}
\newcommand{\Q}{\mathbb Q}
\newcommand{\F}{\mathbb F}
\newcommand{\eps}{\varepsilon}
\begin{document}
% UnsolvedMath ID 30003909; source catalogue code OWR-16411-001
\setcounter{section}{30003908}
\section{Axisymmetry of Spherical Near-Horizon Geometries}\label{30003909}
{\small UnsolvedMath ID 30003909 · Mathematical Physics}
\subsection{Definitions and mathematical statement}

\paragraph{Shared notation.}$\mathbb N=\{1,2,\ldots\}$; the standard number fields have their usual meanings. All additional conventions are stated locally.


\paragraph{Definitions and conventions.}
A degenerate black-hole horizon has zero surface gravity. Its near-horizon limit gives equations for a Riemannian cross-section metric and a one-form on the horizon section. Axisymmetry means there is a nontrivial circle action preserving those data. A spherical section is topologically $S^2$; classification up to diffeomorphism must respect the field equation\textquotesingle s cosmological constant and matter assumptions. A Kerr near-horizon geometry is the vacuum spherical example singled out by the question. In the primary vacuum convention, the unknowns are a smooth metric $g_{AB}$ and one-form $\omega_A$ on $S^2$, satisfying $R_{AB}-(D_A\omega_B+D_B\omega_A+2\omega_A\omega_B)=0$, where $D$ and $R_{AB}$ are the metric covariant derivative and Ricci tensor.

\paragraph{Target retained from the canonical source catalogue.}
\paragraph{Catalogue wording retained for traceability.}
Is every solution of the near-horizon-geometry equations on $S^2$ axisymmetric and therefore induced by a degenerate Kerr metric?

\paragraph{Source scope and qualification.}
The actual vacuum equation was recovered from PDF page 18. The report establishes local rigidity near Kerr and states the global removal-of-axisymmetry conjecture; modern status is not freshly established here. The problem is not enlarged to matter or nonzero cosmological constant.
\subsection{Short English statement}
Must every spherical solution of the indicated near-horizon equations have rotational symmetry and come from extremal Kerr?
\subsection{Sources}
\begin{itemize}
\item[\hypertarget{source:30003909:1}{S1}] ULAM.AI and the UnsolvedMath Contributors, UnsolvedMath, record 30003909 (OWR-16411-001). Catalogue curation under CC BY 4.0; definitions and English exposition added locally.
\url{https://huggingface.co/datasets/ulamai/UnsolvedMath}.
\item[\hypertarget{source:30003909:2}{S2}] Mathematical General Relativity (2018), Oberwolfach Reports, DOI 10.4171/owr/2018/36, pp. 2157--2251. Primary report accessed 2026-10-11; relevant PDF pages 18.
\url{https://doi.org/10.4171/owr/2018/36}.
\item[\hypertarget{source:30003909:3}{S3}] Chrusciel--Szybka--Tod 2018 (catalogue-listed context; not independently consulted in this preparation).
\url{https://doi.org/10.1088/1361-6382/aa90e7}.
\end{itemize}
\label{end:30003909}
\end{document}
